\section{Introduction}
Industrial robot skills are often implemented as a combination of geometric planning, feedback control and task-specific execution contracts. These components can provide reliable approach and transport motions, while sustained object contact remains sensitive to object motion, orientation, friction and tool placement. A locally successful preparation motion need not leave a state from which the next operation finishes stably. Reinforcement learning (RL) offers a way to adjust executable actions using physical task feedback without requiring it to replace inverse kinematics or collision-aware planning.

Our application context is an existing industrial skill architecture, ARES/ARES-R. The present algorithmic experiments, however, are deliberately narrower than that system: a single JAKA arm manipulates a free synthetic object with a research paddle in two consecutive contact stages. They are not claims of tray grasping, production cabinet operation or end-to-end autonomous field execution. Controlled October field records are retained on a separate reference track and are excluded from this public package.

Previous experiments already show that residual PPO can learn these tasks. They also expose an important methodological risk. Aggregate performance on a mixture of tasks can hide a weak result on the more contact-sensitive task, and an auxiliary network or learned switch can update without improving physical outcomes. On unseen geometry, residual PPO improves the aggregate result over a tuned controller, but most of the gain is attributable to P1. Simple waiting at handoff and learned holding do not deliver an independent advantage. These negative findings motivate a bounded investigation rather than further scaling of a weak mechanism.

We focus on the following question: \emph{Can task-relative legal residual parameters, informed by observed execution effects and subsequent task requirements, improve guided contact completion beyond both a strong task-frame controller and ordinary residual RL?} This does not require showing that RL is indispensable. It requires a measurable, repeatable contribution under equivalent information and interaction budgets.

The currently supported contributions are an auditable contact-RL experimental basis and a three-seed historical residual-control comparison. The proposed contribution is an execution-result-aware residual skill mechanism. Its novelty and performance remain to be established through a task-frame baseline, causal action interventions, conditioning ablations and a new held-out evaluation. Merely concatenating context, rotating coordinates or training a predictor is not considered a sufficient algorithmic contribution.

\safewidefigure{Fig1_system_layers.pdf}{Scope of the study within the ARES skill architecture. The contact-RL experiment is implemented; multimodal policy deployment and long-horizon interfaces are not established results.}{fig:system}
